\chapter{Hamiltonian Mechanics and Phase-Space Flow}

Tonight we turn to Hamiltonian mechanics. This is still classical mechanics,
not quantum mechanics, but the change of language will matter enormously when
we reach quantum theory. The immediate purpose is more elementary: we want a
description in which a complete state is a point in phase space and the law of
motion carries that point along a reversible flow.

The route has three stages. First we ask what information a state must contain.
Then a Legendre transformation replaces velocities by canonical momenta.
Finally Hamilton's equations turn one function, the Hamiltonian, into the
velocity field of the entire phase space.

\section{Reversible Laws and Complete States}

Begin with a coin whose state is either heads or tails. Imagine that time
advances in discrete steps. One possible rule leaves the coin unchanged:
heads goes to heads and tails goes to tails. Another rule flips it at every
step. In either case the present determines not only the next state but also
the preceding one. The evolution can be run forward or backward without
ambiguity.

For a system with many possible states, represent each state by a point and
draw an arrow from every point to its successor. A reversible rule permits
exactly one arrow into and one arrow out of every point. For a finite state
space, the arrows must therefore decompose into cycles. If the cycles are
marked by different labels, such as \(+1\) and \(-1\), the label is conserved:
a trajectory cannot jump from one cycle to another.

\begin{figure}[tbp]
\centering
\includegraphics[width=0.82\textwidth]{lecture_06_figure_01.png}
\caption{Discrete evolution at 00:08:05. The central cycles belong to
separate sectors, while the map at the left illustrates how several distinct
histories can merge when backward evolution is not unique.}
\label{fig:l06-discrete-flow}
\end{figure}

Now contrast this with a many-to-one rule. Several states may have the same
successor. Forward prediction remains possible, but retrodiction does not:
after reaching the common image, we cannot infer which predecessor was
occupied. Distinguishable histories have been erased.

The same idea can be pictured by covering the state space with red dots and
letting every dot move. Under a reversible rule no dot disappears, splits, or
merges with another. In a continuous Hamiltonian system something stronger is
true: phase-space volume is preserved. The points move as an incompressible
fluid. Reversibility by itself does not prove incompressibility for every
conceivable continuous flow; the canonical structure of Hamilton's equations
does. We shall see the local reason for that structure in the crossed
derivatives of those equations.

The phrase \emph{phase space} will mean the space of complete instantaneous
states. A point in it must carry enough information for the equations of
motion to determine a unique trajectory. If the microscopic law is reversible,
the same point also determines the past.

\section{Why Position Is Not a State}

Suppose a particle is observed at a position \(q\). Where will it be one
second later? Position alone does not answer the question. The particle may be
moving in any direction and at any speed. A complete state must include a
second set of data,
\begin{equation}
 (q_i,\dot q_i)
 \qquad\text{or, after a change of variables,}\qquad
 (q_i,p_i).
 \label{eq:l06-complete-state}
\end{equation}
The coordinates \(q_i\) form configuration space. The pairs \((q_i,p_i)\)
form phase space.

This distinction is already visible in Newton's law. For \(N\) Cartesian
coordinates,
\begin{equation}
 m\ddot x_i=F_i,
 \qquad i=1,\ldots,N,
 \label{eq:l06-newton-second-order}
\end{equation}
is a set of \(N\) second-order equations. Define
\begin{equation}
 p_i=m\dot x_i.
 \label{eq:l06-cartesian-momentum}
\end{equation}
The same physics becomes \(2N\) first-order equations,
\begin{equation}
 \boxed{
 \dot x_i=\frac{p_i}{m},
 \qquad
 \dot p_i=F_i.}
 \label{eq:l06-newton-first-order}
\end{equation}

\begin{figure}[tbp]
\centering
\includegraphics[width=0.78\textwidth]{lecture_06_figure_02.png}
\caption{The Newton equation rewritten as two first-order equations at
00:18:48. No new physics has been added; the complete state has been made
explicit.}
\label{fig:l06-first-order-newton}
\end{figure}

This rewriting is not yet the full Hamiltonian construction. Any ordinary
second-order equation can be made first order by introducing an extra
variable. Hamiltonian mechanics demands the more special pairing of
coordinates and canonical momenta, with both first-order equations generated
by derivatives of a single function.

\begin{classroomqa}
\audiencequestion{Why reformulate \(F=ma\) if it already predicts the
motion?}
\lecturerresponse{The Lagrangian formulation freed us from Cartesian
coordinates. Motion on a plane may be described by \(x\) and \(y\), while a
pendulum or a particle on a sphere is more naturally described by angles.
Hamilton's formulation reveals an additional structure: mechanics becomes a
flow on a space of canonical pairs. The reformulation is useful in classical
mechanics and becomes indispensable in quantum mechanics. Historical motives
are harder to reconstruct than the mathematics, so the important test is what
the construction accomplishes.}
\end{classroomqa}

Imagine placing one identical system at every phase-space point. Each copy
has an instantaneous velocity \((\dot q_i,\dot p_i)\), so the equations define
a vector field on phase space. Following that field carries the whole cloud
of systems forward. Our task is to find the function that generates this
flow.

\section{The Legendre Transformation}

The required mathematical device is a Legendre transformation. Start with
one variable \(V\) and a differentiable function \(L(V)\). Define the
conjugate variable
\begin{equation}
 P=\frac{dL}{dV}.
 \label{eq:l06-legendre-p}
\end{equation}
Assume, at least in the region of interest, that this relation can be inverted
to give \(V=V(P)\). Geometrically the graph of \(P(V)\) must not double back
in a way that assigns several relevant values of \(V\) to the same \(P\).

With convenient additive constants one may write
\begin{equation}
 L(V)=\int_0^V P(V')\,dV',
 \qquad
 H(P)=\int_0^P V(P')\,dP'.
 \label{eq:l06-legendre-areas}
\end{equation}
For a monotone graph through the chosen origin, these are complementary areas
inside a rectangle of area \(PV\):
\begin{equation}
 H(P)+L(V)=PV.
 \label{eq:l06-complementary-area}
\end{equation}
The invariant definition, independent of the choice of additive constants,
is
\begin{equation}
 \boxed{H(P)=PV-L(V),}
 \qquad P=L'(V),
 \label{eq:l06-legendre-definition}
\end{equation}
where \(V\) on the right is to be expressed as a function of \(P\).

\begin{figure}[tbp]
\centering
\includegraphics[width=0.72\textwidth]{lecture_06_figure_03.png}
\caption{The complementary-area construction at 00:32:05. The area labeled
\(L(V)\) lies under \(P(V)\); the remaining part of the \(PV\) rectangle is
the transformed function \(H(P)\).}
\label{fig:l06-legendre-area}
\end{figure}

The derivative of the transformed function follows by varying
Eq.~\eqref{eq:l06-legendre-definition}:
\begin{align}
 \delta H
 &=P\,\delta V+V\,\delta P
   -\frac{dL}{dV}\,\delta V \notag\\
 &=P\,\delta V+V\,\delta P-P\,\delta V
 =V\,\delta P.
 \label{eq:l06-legendre-cancellation}
\end{align}
Therefore
\begin{equation}
 \boxed{\frac{dH}{dP}=V.}
 \label{eq:l06-legendre-inverse}
\end{equation}
Differentiation created \(P\) from \(L(V)\); after the transform,
differentiation creates \(V\) from \(H(P)\). This reciprocity is the point of
the construction.

\begin{figure}[tbp]
\centering
\includegraphics[width=0.82\textwidth]{lecture_06_figure_04.png}
\caption{The cancellation in the variation of \(H=PV-L\) at 00:36:30.
The boxed terms cancel because \(dL/dV=P\), leaving
\(\delta H=V\,\delta P\).}
\label{fig:l06-legendre-variation}
\end{figure}

For several variables, let \(L=L(v_1,\ldots,v_N)\) and define
\begin{equation}
 p_i=\frac{\partial L}{\partial v_i}.
 \label{eq:l06-multivariable-p}
\end{equation}
If these equations can be inverted for all the \(v_i\), the transform is
\begin{equation}
 H(p_1,\ldots,p_N)=\sum_i p_i v_i-L,
 \label{eq:l06-multivariable-H}
\end{equation}
and the same cancellation gives
\begin{equation}
 \frac{\partial H}{\partial p_i}=v_i.
 \label{eq:l06-multivariable-inverse}
\end{equation}
Legendre transformations appear repeatedly in thermodynamics, where changing
which variables are held fixed produces the familiar thermodynamic
potentials. Here we use the same mathematics to exchange velocities for
momenta.

\subsection{Coordinates Go Along for the Ride}

In mechanics the Lagrangian is not merely \(L(v)\), but
\begin{equation}
 L=L(q,v,t),
 \qquad v_i=\dot q_i.
 \label{eq:l06-L-qvt}
\end{equation}
The transformation is performed only in the velocity variables. During it,
the coordinates \(q_i\) and time \(t\) are held fixed as parameters:
\begin{equation}
 p_i=\frac{\partial L}{\partial v_i},
 \qquad
 H(q,p,t)=\sum_i p_i v_i-L(q,v,t).
 \label{eq:l06-mechanical-transform}
\end{equation}
The inverted velocities \(v_i=v_i(q,p,t)\) are then substituted on the
right, so \(H\) is genuinely a function of \(q,p,t\).

\begin{classroomqa}
\audiencequestion{May the velocity determined by the momentum also depend
on position?}
\lecturerresponse{Yes. The Legendre transformation holds \(q\) fixed, but
the relation \(p_i=\partial L/\partial v_i\) may contain \(q\). A particle
in curvilinear coordinates is a standard example: the coefficients in the
kinetic energy depend on position, so solving for \(v_i\) generally produces
\(v_i(q,p)\).}
\end{classroomqa}

Local invertibility is controlled by the velocity Hessian
\begin{equation}
 W_{ij}=\frac{\partial^2L}{\partial v_i\partial v_j}.
 \label{eq:l06-velocity-hessian}
\end{equation}
For ordinary unconstrained mechanics, \(\det W\neq0\) allows the velocities
to be solved locally in terms of the momenta. A singular Hessian does not mean
that nature is inconsistent; it signals a constrained system requiring a
more careful Hamiltonian treatment.

\section{Hamilton's Equations}

Now vary the mechanical transform while allowing both \(q_i\) and \(p_i\) to
change. At fixed time,
\begin{align}
 \delta H
 &=
 \sum_i\left(p_i\,\delta v_i+v_i\,\delta p_i\right)
 -\sum_i\frac{\partial L}{\partial q_i}\,\delta q_i
 -\sum_i\frac{\partial L}{\partial v_i}\,\delta v_i \notag\\
 &=
 \sum_i v_i\,\delta p_i
 -\sum_i\frac{\partial L}{\partial q_i}\,\delta q_i.
 \label{eq:l06-full-variation}
\end{align}
The two terms proportional to \(\delta v_i\) cancel because
\(p_i=\partial L/\partial v_i\). Reading off the coefficients gives
\begin{equation}
 \frac{\partial H}{\partial p_i}=v_i=\dot q_i,
 \qquad
 \frac{\partial H}{\partial q_i}
 =-\frac{\partial L}{\partial q_i}.
 \label{eq:l06-H-partials}
\end{equation}

The Euler--Lagrange equations say
\begin{equation}
 \frac{d}{dt}\frac{\partial L}{\partial\dot q_i}
 =\frac{\partial L}{\partial q_i}.
 \label{eq:l06-EL-again}
\end{equation}
Since the derivative on the left is \(\dot p_i\), we arrive at
\begin{equation}
 \boxed{
 \dot q_i=\frac{\partial H}{\partial p_i},
 \qquad
 \dot p_i=-\frac{\partial H}{\partial q_i}.}
 \label{eq:l06-hamilton-equations}
\end{equation}

\begin{figure}[tbp]
\centering
\includegraphics[width=0.82\textwidth]{lecture_06_figure_05.png}
\caption{The completed pair of Hamilton equations at 00:52:40. The crossed
derivatives and the single minus sign encode the canonical phase-space
structure.}
\label{fig:l06-hamilton-board}
\end{figure}

The equations are first order and nearly symmetric in \(q_i\) and \(p_i\).
Their derivatives are crossed: the rate of change of \(q_i\) comes from a
momentum derivative, while the rate of change of \(p_i\) comes from a
coordinate derivative. The minus sign is essential. It makes the phase-space
velocity field
\begin{equation}
 \bm{u}_H=
 \left(
 \frac{\partial H}{\partial p_1},\ldots,
 \frac{\partial H}{\partial p_N},
 -\frac{\partial H}{\partial q_1},\ldots,
 -\frac{\partial H}{\partial q_N}
 \right)
 \label{eq:l06-phase-velocity}
\end{equation}
divergence-free:
\begin{equation}
 \nabla_{q,p}\!\cdot\bm{u}_H
 =\sum_i
 \left(
 \frac{\partial^2H}{\partial q_i\partial p_i}
 -\frac{\partial^2H}{\partial p_i\partial q_i}
 \right)=0.
 \label{eq:l06-liouville-local}
\end{equation}
Thus the red-dot fluid is locally incompressible. This is the elementary
coordinate form of Liouville's theorem.

\begin{classroomqa}
\audiencequestion{Is Hamiltonian mechanics merely the trick of replacing
each second-order equation by two first-order equations?}
\lecturerresponse{No. Doubling the variables is always possible, but an
arbitrary doubled system need not have the crossed-gradient form in
Eq.~\eqref{eq:l06-hamilton-equations}. The canonical form, with one
Hamiltonian generating both equations, is the stronger statement.}
\end{classroomqa}

\begin{classroomqa}
\audiencequestion{What happens if the relation between velocity and
momentum cannot be inverted?}
\lecturerresponse{For the regular systems being discussed, the ordinary
Legendre transform then fails and \(H(q,p)\) cannot be constructed by this
recipe. Such a model may be pathological, but singular Lagrangians also arise
in legitimate constrained and gauge systems. They require additional
constraint equations rather than the naive transform.}
\end{classroomqa}

\begin{classroomqa}
\audiencequestion{Does friction, including terminal velocity, violate the
Hamiltonian picture?}
\lecturerresponse{A friction law for a reduced set of variables is generally
not Hamiltonian: phase-space volumes contract and mechanical information is
discarded. At the microscopic level, however, the object and its environment
may form a larger Hamiltonian system. Energy and information then flow into
unobserved molecular, acoustic, or radiative degrees of freedom.}
\end{classroomqa}

\begin{classroomqa}
\audiencequestion{Where, exactly, did the Legendre transformation enter the
mechanics?}
\lecturerresponse{It is the passage from \(L(q,v)\) to
\(H(q,p)=\sum_i p_i v_i-L\), with
\(p_i=\partial L/\partial v_i\) and the velocities eliminated in favor of
the momenta. The geometric area picture explains the reciprocal derivatives;
the cancellation in Eq.~\eqref{eq:l06-full-variation} is what turns that
reciprocity into Hamilton's equations.}
\end{classroomqa}

\begin{classroomqa}
\audiencequestion{Can a Legendre transform fail because the velocity becomes
singular or several velocities correspond to one momentum?}
\lecturerresponse{Yes. Global one-to-one behavior can fail even if a local
branch exists, and the Hessian can become singular at special points. Before
using the ordinary Hamiltonian construction one must specify a region where
the momentum--velocity relation is regular, or adopt the constrained
formalism appropriate to the system.}
\end{classroomqa}

\begin{classroomqa}
\audiencequestion{Do infinitely many degrees of freedom, as in a field,
automatically produce friction?}
\lecturerresponse{No. A classical field can have infinitely many canonical
pairs and still evolve Hamiltonianly. Apparent damping arises when energy is
carried away into modes that are omitted from the reduced description. An
accelerating charge, for example, may lose mechanical energy to radiation,
while the combined particle-plus-field system retains the full accounting.}
\end{classroomqa}

\section{The Ordinary Particle}

For one particle in a potential \(U(x)\),
\begin{equation}
 L(x,\dot x)=\frac12m\dot x^2-U(x).
 \label{eq:l06-particle-L}
\end{equation}
The canonical momentum is
\begin{equation}
 p=\frac{\partial L}{\partial\dot x}=m\dot x,
 \qquad
 \dot x=\frac{p}{m}.
 \label{eq:l06-particle-p}
\end{equation}
Its Legendre transform is
\begin{align}
 H
 &=p\dot x-L \notag\\
 &=p\frac{p}{m}
   -\left(\frac12m\frac{p^2}{m^2}-U(x)\right) \notag\\
 &=\boxed{\frac{p^2}{2m}+U(x)}.
 \label{eq:l06-particle-H}
\end{align}
For this familiar Lagrangian, the Hamiltonian is the total energy \(T+U\).

\begin{figure}[tbp]
\centering
\includegraphics[width=0.76\textwidth]{lecture_06_figure_06.png}
\caption{The ordinary-particle Hamiltonian at 01:12:35. The velocity has
been eliminated, leaving \(H=p^2/(2m)+U(x)\).}
\label{fig:l06-particle-board}
\end{figure}

The first Hamilton equation gives
\begin{equation}
 \dot x=\frac{\partial H}{\partial p}=\frac{p}{m},
 \label{eq:l06-first-H-particle}
\end{equation}
which is the momentum--velocity relation. The second gives
\begin{equation}
 \dot p=-\frac{\partial H}{\partial x}
 =-\frac{dU}{dx}.
 \label{eq:l06-second-H-particle}
\end{equation}
Since \(p=m\dot x\), this is Newton's equation
\begin{equation}
 m\ddot x=-\frac{dU}{dx}.
\end{equation}
The reformulation has changed the organization of the equations, not their
physical predictions.

\section{Energy and the Geometry of Its Contours}

Suppose \(H(q,p)\) has no explicit time dependence. Along a trajectory,
\begin{align}
 \frac{dH}{dt}
 &=\sum_i\left(
 \frac{\partial H}{\partial q_i}\dot q_i
 +\frac{\partial H}{\partial p_i}\dot p_i
 \right) \notag\\
 &=\sum_i\left(
 \frac{\partial H}{\partial q_i}
 \frac{\partial H}{\partial p_i}
 -\frac{\partial H}{\partial p_i}
 \frac{\partial H}{\partial q_i}
 \right)=0.
 \label{eq:l06-energy-cancellation}
\end{align}
The cancellation is pairwise, so
\begin{equation}
 \boxed{H=\text{constant along the motion}.}
 \label{eq:l06-H-conserved}
\end{equation}

\begin{figure}[tbp]
\centering
\includegraphics[width=0.82\textwidth]{lecture_06_figure_07.png}
\caption{Energy conservation at 01:16:10. Substituting Hamilton's equations
into \(dH/dt\) makes the two products cancel term by term.}
\label{fig:l06-energy-board}
\end{figure}

\begin{classroomqa}
\audiencequestion{Is the converse true? Does energy conservation by itself
imply Hamiltonian dynamics?}
\lecturerresponse{No. A time-independent Hamiltonian evolving by Hamilton's
equations is conserved, but a dynamical rule may preserve some function
called energy without possessing canonical Hamiltonian structure. Energy
conservation is one consequence, not a complete characterization, of
Hamiltonian flow.}
\end{classroomqa}

For one degree of freedom, constant values of \(H(q,p)\) are curves in the
\(q\)-\(p\) plane. The phase point must move along the curve on which it
starts. For a harmonic oscillator,
\begin{equation}
 H=\frac{p^2}{2m}+\frac12m\omega^2q^2,
 \label{eq:l06-oscillator-H}
\end{equation}
the curves are ellipses, or circles after rescaling the axes. The origin is a
fixed point. Every nonzero bounded orbit closes around it.

\begin{figure}[tbp]
\centering
\begin{tikzpicture}[x=1.05cm,y=1.05cm,>=Stealth]
  \begin{scope}
    \draw[->] (-3.0,0) -- (3.1,0) node[right] {\(q\)};
    \draw[->] (0,-2.25) -- (0,2.35) node[above] {\(p\)};
    \draw[blue!60!black,thick] (0,0) ellipse (1.15 and 0.75);
    \draw[blue!60!black,thick] (0,0) ellipse (2.05 and 1.35);
    \draw[blue!60!black,thick] (0,0) ellipse (2.75 and 1.85);
    \draw[->,very thick,orange!75!black]
      (2.05,0) arc[start angle=0,end angle=36,x radius=2.05,y radius=1.35];
    \fill (0,0) circle (1.7pt);
    \node at (0,-2.65) {harmonic oscillator};
  \end{scope}
  \begin{scope}[xshift=8.0cm]
    \draw[->] (-3.0,0) -- (3.1,0) node[right] {\(q\)};
    \draw[->] (0,-2.25) -- (0,2.35) node[above] {\(p\)};
    \draw[gray!70,thick,samples=140,domain=-3:3]
      plot (\x,{0.38*cos(90*\x)});
    \draw[green!45!black,thick]
      plot[smooth] coordinates {(-2.9,-1.65) (-2.1,-1.25)
      (-1.35,-0.55) (-0.75,-0.08) (-0.2,0.24) (0.45,0.58)
      (1.2,0.94) (2.0,1.25) (2.9,1.56)};
    \draw[red!65!black,thick]
      plot[smooth] coordinates {(-1.42,0) (-1.22,0.56) (-0.78,0.94)
      (0,1.08) (0.78,0.94) (1.22,0.56) (1.42,0)
      (1.22,-0.56) (0.78,-0.94) (0,-1.08)
      (-0.78,-0.94) (-1.22,-0.56) (-1.42,0)};
    \node[red!65!black] at (0,-1.45) {bounded};
    \node[green!45!black] at (1.8,1.75) {running};
    \node at (0,-2.65) {periodic potential};
  \end{scope}
\end{tikzpicture}
\caption{Typeset phase portraits. Small energy in a stable well produces a
closed orbit. Above a periodic potential barrier, the coordinate can keep
running and the constant-energy curve is open.}
\label{fig:l06-phase-portraits}
\end{figure}

Closed and open curves can coexist in one phase space. A pendulum below the
separatrix oscillates and traces a closed curve. Above the barrier it rotates
continuously, producing an open branch when the angle is unwrapped. The same
distinction appears for a particle moving through a periodic washboard
potential.

\begin{classroomqa}
\audiencequestion{If the phase-space fluid is incompressible, how can
neighboring energy contours bunch together and spread apart?}
\lecturerresponse{Think of an incompressible fluid passing through a narrow
channel. Where the channel narrows, the fluid moves faster; where it widens,
the fluid moves more slowly. Closely spaced contours mean that the gradient
of \(H\) is large, and Hamilton's equations therefore give a large
phase-space speed. The density need not change.}
\end{classroomqa}

\begin{classroomqa}
\audiencequestion{May closed and open constant-energy curves occur in the
same system?}
\lecturerresponse{Yes. A periodic potential provides the simplest example.
Below a barrier, the particle is trapped in one well and the orbit is closed.
Above it, the particle passes from well to well and the corresponding branch
is open. The pendulum has the same oscillating and rotating regimes.}
\end{classroomqa}

\subsection{Momentum Need Not Be Mass Times Velocity}

The relation \(p=m\dot q\) belongs to the ordinary quadratic kinetic energy;
it is not the definition of canonical momentum in general. Hamilton's first
equation always says
\begin{equation}
 \dot q=\frac{\partial H}{\partial p}.
 \label{eq:l06-general-p-v}
\end{equation}
For the illustrative Hamiltonian written in the lecture,
\begin{equation}
 H=p^3+q^3,
 \label{eq:l06-cubic-H}
\end{equation}
one finds
\begin{equation}
 \dot q=3p^2,
 \qquad
 \dot p=-3q^2.
 \label{eq:l06-cubic-flow}
\end{equation}

\begin{figure}[tbp]
\centering
\includegraphics[width=0.76\textwidth]{lecture_06_figure_08.png}
\caption{A deliberately nonstandard Hamiltonian at 01:25:50. It shows
directly that the canonical momentum need not be proportional to velocity.}
\label{fig:l06-nonlinear-board}
\end{figure}

\begin{classroomqa}
\audiencequestion{Is canonical momentum always a constant multiple of
velocity?}
\lecturerresponse{No. It is defined by
\(p_i=\partial L/\partial\dot q_i\), or equivalently related to velocity by
\(\dot q_i=\partial H/\partial p_i\). That relation may be nonlinear and
may depend on the coordinates. The cubic example makes the nonlinearity
obvious, although its momentum--velocity map is not globally one-to-one and
therefore does not come from one globally regular ordinary Lagrangian.}
\end{classroomqa}

\begin{classroomqa}
\audiencequestion{If microscopic Hamiltonian motion conserves energy, does
that make every machine a perpetual-motion machine?}
\lecturerresponse{No. Total energy conservation does not mean that all
energy remains available as organized mechanical work. Energy can spread
among enormous numbers of environmental degrees of freedom and become heat.
Understanding that irreversibility and the distinction between total and
usable energy requires statistical mechanics, entropy, and often chaotic
dynamics.}
\end{classroomqa}

\section{Poisson Brackets and General Conservation Laws}

Energy is not the only quantity that may be conserved. Let \(A(q,p)\) be any
function on phase space. Angular momentum in a plane is one example:
\begin{equation}
 L_z=xp_y-yp_x.
 \label{eq:l06-Lz}
\end{equation}
For an observable with no explicit time dependence, the chain rule and
Hamilton's equations give
\begin{align}
 \frac{dA}{dt}
 &=\sum_i\left(
 \frac{\partial A}{\partial q_i}\dot q_i+
 \frac{\partial A}{\partial p_i}\dot p_i
 \right) \notag\\
 &=\sum_i\left(
 \frac{\partial A}{\partial q_i}
 \frac{\partial H}{\partial p_i}
 -
 \frac{\partial A}{\partial p_i}
 \frac{\partial H}{\partial q_i}
 \right).
 \label{eq:l06-A-dot}
\end{align}

This combination is the Poisson bracket. With the convention used here,
\begin{equation}
 \boxed{
 \{A,B\}
 =\sum_i\left(
 \frac{\partial A}{\partial q_i}
 \frac{\partial B}{\partial p_i}
 -
 \frac{\partial A}{\partial p_i}
 \frac{\partial B}{\partial q_i}
 \right).}
 \label{eq:l06-poisson-definition}
\end{equation}
Consequently,
\begin{equation}
 \boxed{\dot A=\{A,H\}.}
 \label{eq:l06-poisson-evolution}
\end{equation}
The Hamiltonian is the generator of time evolution.

\begin{figure}[tbp]
\centering
\includegraphics[width=0.88\textwidth]{lecture_06_figure_09.png}
\caption{The chain rule and Poisson-bracket definition at 01:35:20. The
bracket packages the crossed derivatives that already appeared in Hamilton's
equations.}
\label{fig:l06-poisson-board}
\end{figure}

Taking \(A=q_i\) or \(A=p_i\) recovers Hamilton's equations:
\begin{equation}
 \{q_i,H\}=\frac{\partial H}{\partial p_i}=\dot q_i,
 \qquad
 \{p_i,H\}=-\frac{\partial H}{\partial q_i}=\dot p_i.
 \label{eq:l06-bracket-Hamilton}
\end{equation}
The fundamental brackets are
\begin{equation}
 \{q_i,q_j\}=0,
 \qquad
 \{p_i,p_j\}=0,
 \qquad
 \{q_i,p_j\}=\delta_{ij}.
 \label{eq:l06-canonical-brackets}
\end{equation}

A quantity \(A(q,p)\) is conserved precisely when
\begin{equation}
 \boxed{\{A,H\}=0.}
 \label{eq:l06-conservation-bracket}
\end{equation}
For a rotationally invariant Hamiltonian, direct differentiation verifies
\(\{L_z,H\}=0\). Later we will connect this calculation to symmetry; for now
the bracket gives the mathematical test.

\begin{classroomqa}
\audiencequestion{How could Poisson have discovered such a combination of
derivatives?}
\lecturerresponse{Historical reconstruction is uncertain, but the
combination is not arbitrary once Hamilton's equations are written. Apply
the chain rule to an arbitrary \(A(q,p)\), substitute the two canonical
equations, and the bracket appears. Its importance is then confirmed by the
algebraic properties and conservation laws it organizes.}
\end{classroomqa}

\section{Explicit Time Dependence}

So far time entered only implicitly through the moving variables \(q(t)\)
and \(p(t)\). An external apparatus can also make the Hamiltonian depend on time
explicitly. A potential \(U(q,t)\), for example, may be produced by a source
whose prescribed position changes while the source itself is not included
among the dynamical variables.

For a general observable \(A(q,p,t)\),
\begin{equation}
 \frac{dA}{dt}=\{A,H\}+\frac{\partial A}{\partial t}.
 \label{eq:l06-explicit-A}
\end{equation}
Setting \(A=H\) gives
\begin{equation}
 \frac{dH}{dt}
 =\{H,H\}+\frac{\partial H}{\partial t}
 =\boxed{\frac{\partial H}{\partial t}},
 \label{eq:l06-explicit-H}
\end{equation}
because every function has zero Poisson bracket with itself. Thus a
time-independent Hamiltonian is conserved, while explicit time dependence
measures the energy supplied to or removed from the chosen subsystem.

\begin{figure}[tbp]
\centering
\includegraphics[width=0.88\textwidth]{lecture_06_figure_10.png}
\caption{Explicit and implicit time dependence separated at 01:45:55.
Hamilton's equations cancel the terms involving \(\dot q_i\) and
\(\dot p_i\), leaving \(dH/dt=\partial H/\partial t\).}
\label{fig:l06-explicit-time-board}
\end{figure}

The same conclusion follows directly from the time-dependent Legendre
transform:
\begin{equation}
 dH
 =\sum_i v_i\,dp_i
 -\sum_i\dot p_i\,dq_i
 -\frac{\partial L}{\partial t}\,dt,
 \label{eq:l06-dH-time}
\end{equation}
so
\begin{equation}
 \frac{\partial H}{\partial t}
 =-\frac{\partial L}{\partial t}
 \label{eq:l06-time-transform}
\end{equation}
when the canonical variables are held fixed on their respective sides of the
transform.

\begin{classroomqa}
\audiencequestion{What is the difference between implicit and explicit time
dependence?}
\lecturerresponse{The value \(H(q(t),p(t))\) changes its arguments as the
system moves; that is implicit dependence. Explicit dependence remains even
when \(q\) and \(p\) are held fixed, as in \(U(q,t)\) produced by a moving
external source. The canonical terms cancel in \(dH/dt\); only the explicit
partial derivative remains.}
\end{classroomqa}

\begin{classroomqa}
\audiencequestion{Why do some books reverse the order or sign of the
Poisson bracket? Which slot represents momentum?}
\lecturerresponse{The names \(A\) and \(B\) label whole phase-space
functions; neither argument slot is the momentum slot. The canonical
variables \(q_i,p_i\) occur inside the definition. Since
\(\{A,B\}=-\{B,A\}\), reversing the order reverses the sign. Either convention
is consistent if Hamilton's evolution equation is adjusted with it. Here we
use \(\dot A=\{A,H\}\).}
\end{classroomqa}

\begin{classroomqa}
\audiencequestion{Is it the Poisson bracket that appears everywhere outside
mechanics?}
\lecturerresponse{The broad remark at this point concerns the Legendre
transformation, which recurs throughout thermodynamics and other variational
theories. Poisson brackets are fundamental in Hamiltonian mechanics and lead
directly toward quantum commutators, but they are a different piece of the
construction.}
\end{classroomqa}

We began with a reversible map of abstract states and ended with a precise
continuous law. The Legendre transformation changes variables from
\((q,\dot q)\) to \((q,p)\); Hamilton's equations generate an
incompressible phase-space flow; constant \(H\) confines that flow to energy
surfaces; and the Poisson bracket tells how every observable changes. These
are not separate tricks. They are four views of the same canonical
structure.
